\chapter{The conjecture, special cases, and status}
\label{ch:conjecture-and-examples}

\section{The property for one group}
\begin{definition}\label{def:baum-connes-property}
\uses{def:assembly-map}
\lean{BC4lean.Assembly.BaumConnesProperty,BC4lean.Assembly.baumConnesProperty_iff_even_and_odd}
\leanok
Write $\mathrm{BC}(\Gamma)$ for the proposition that $\mu_{\Gamma,i}$
is an isomorphism for both $i=0$ and $i=1$.
\end{definition}

\section{The general conjecture}
\begin{conjecture}\label{conj:baum-connes}
\uses{def:baum-connes-property}
\lean{BC4lean.Assembly.BaumConnesConjecture}
\notready
Every countable discrete group $\Gamma$ satisfies $\mathrm{BC}(\Gamma)$.
\end{conjecture}

The linked Lean declaration encodes this quantified proposition in each
fixed Lean universe using the actual canonical reduced assembly maps.
The proposition remains unproved; its definition supplies no witness and
introduces no axiom asserting it.

\section{Calibration examples and later specializations}
For the literal group $\mathrm{Unit}:\mathrm{Type}$, the proof below
constructs an additive equivalence and identifies it with the actual canonical
assembly map in each degree. For a general finite group, the result below
reduces the actual assembly map to its one-point cutoff/product/descent
formula; it does not assert bijectivity. Numerical $K$-group computations
and further assembly isomorphisms for finite groups and $\mathbb Z$ remain
later developments. Treat $\mathrm{SL}(3,\mathbb Z)$ as a later specialization;
no result for it follows merely from the constructions here.


\section{Constructions, comparisons, and the scope of the conjecture}
The linked declaration \texttt{BaumConnesConjecture} encodes a proposition
in each fixed Lean universe. Defining that proposition supplies no witness
that it holds. The coefficient-free reduced conjecture above remains an
unproved target; the lemmas below distinguish its exact logical meaning,
comparison results, and a separate proof for the literal Unit group.

\begin{proposition}\label{prop:baum-connes-model-invariance}
\uses{def:baum-connes-property,prop:assembly-compatibility}
\lean{BC4lean.Assembly.assemblyHom_injective_iff_homotopyEquiv,BC4lean.Assembly.assemblyHom_surjective_iff_homotopyEquiv,BC4lean.Assembly.assemblyHom_bijective_iff_homotopyEquiv,BC4lean.Assembly.modelAssemblyProperty_iff_baumConnesProperty}
\leanok
For the indicated proper locally compact Hausdorff second-countable spaces,
an actual equivariant homotopy equivalence preserves injectivity,
surjectivity and bijectivity of the constructed assembly map in each degree.
For genuine universal proper CW models satisfying the subgroup fixed-point
criterion, the two-degree property agrees with the property of the canonical
geometric telescope. Universal-model comparison retains the common group
and space universe required by the existing API; homotopy comparison permits
an independent group universe.
\end{proposition}
\begin{proof}
\leanok
Use the constructed topological additive equivalence and the proved equality
of actual assembly values under its comparison map. Transport each of
injectivity and surjectivity through that equivalence. The universal-model
comparison is supplied by the proved CW fixed-point uniqueness theorem.
No assembly bijectivity hypothesis is introduced by the comparison theorem.
\end{proof}

\begin{proposition}\label{prop:baum-connes-cofinal-invariance}
\uses{prop:baum-connes-model-invariance,prop:assembly-cofinal-index}
\lean{BC4lean.Assembly.indexedAssemblyHom_injective_iff_cofinalIndex,BC4lean.Assembly.indexedAssemblyHom_surjective_iff_cofinalIndex,BC4lean.Assembly.indexedAssemblyHom_bijective_iff_cofinalIndex,BC4lean.Assembly.indexedAssemblyProperty_iff_cofinalComparison}
\leanok
Assembly injectivity, surjectivity and bijectivity are unchanged by a monotone
restriction $\phi$ of the cocompact-piece indexing system satisfying the full
cofinal condition $\forall S\,\exists j,\ S\subseteq\phi(j)$. The two-degree
properties agree for any two such cofinal restrictions. Pointwise covering
alone is not the cofinal hypothesis of this result.
\end{proposition}
\begin{proof}
\leanok
Use the actual cofinal colimit equivalence and the previously proved equality
of assembly values under that equivalence. Compare two restrictions by their
common full cocompact-piece system.
\end{proof}

\begin{proposition}\label{prop:assembly-piece-bijectivity-criteria}
\uses{def:assembly-map,prop:piece-assembly-transition}
\lean{BC4lean.Assembly.assemblyHom_surjective_iff_pieceImages,BC4lean.Assembly.assemblyHom_injective_iff_commonTransition,BC4lean.Assembly.assemblyHom_bijective_iff_pieceCriteria}
\leanok
The actual assembly map is surjective exactly when every analytic class is
the actual assembly value of a class on some cocompact piece. It is injective
exactly when any two piece classes with equal assembly values become equal
after transition to one common larger cocompact piece. Its bijectivity is the
conjunction of these two conditions. The equivalences assert neither
condition for an arbitrary group.
\end{proposition}
\begin{proof}
\leanok
Every colimit class has an actual piece representative; assembly restricts
to the genuine piece formula. Equality in this actual colimit is precisely
equality after passage to one common larger piece. Apply these representative
and equality characterizations in both directions.
\end{proof}

\begin{proposition}\label{prop:trivial-group-analytic-calibration}
\uses{prop:reduced-group-cstar-trivial,def:operator-k0,def:operator-k1}
\lean{BC4lean.Chapter8Examples.trivialAnalyticK0Equiv,BC4lean.Chapter8Examples.trivialAnalyticK1Equiv,BC4lean.Chapter8Examples.trivialAnalyticK0Equiv_projection,BC4lean.Chapter8Examples.trivialAnalyticK1Equiv_unitary}
\leanok
For a one-element group, the actual reduced group C$^*$-algebra scalar
identification induces additive equivalences between its actual reduced
$K_0$ and $K_1$ targets and the corresponding operator $K$-groups of
$\mathbb C$. The maps retain the actual entrywise projection and unitary
class formulae. This statement does not compute $K_0(\mathbb C)$ as
$\mathbb Z$, compute $K_1(\mathbb C)$ as zero, or identify assembly.
\end{proposition}
\begin{proof}
\leanok
Apply the proved functorial equivalences for unital projection $K_0$ and
stable-unitary $K_1$ to the actual scalar star-algebra equivalence, retaining
the spectral orders used by the reduced analytic targets.
\end{proof}

\begin{proposition}\label{prop:compact-model-colimit-calibration}
\uses{def:assembly-map,prop:piece-assembly-transition}
\lean{BC4lean.Chapter8Examples.wholeCompactPiece,BC4lean.Chapter8Examples.terminalColimitEquiv,BC4lean.Chapter8Examples.compactSpaceGroupEquiv,BC4lean.Chapter8Examples.assemblyHom_eq_compactSpaceAssembly}
\leanok
On an indicated compact ambient space, the actual topological colimit is
additively equivalent to the actual full equivariant KK group of the ambient
space. On a proper such space, this equivalence identifies colimit assembly
with the actual ambient-space cutoff/product/descent formula.
\end{proposition}
\begin{proof}
\leanok
The whole compact space is a closed invariant cocompact piece and is greatest
in the piece order. Construct colimit evaluation by transitions into this
piece and use its canonical inclusion for the inverse; prove both inverse
laws. The actual all-space subtype is equivariantly homeomorphic to the
ambient space. Its genuine proper maps induce mutually inverse native KK
maps. Transition compatibility and actual proper assembly naturality give
the assembly equality.
\end{proof}

\begin{proposition}\label{prop:finite-group-point-reduction}
\uses{prop:compact-model-colimit-calibration,prop:assembly-compatibility}
\lean{BC4lean.Chapter8Examples.FiniteGroupPoint,BC4lean.Chapter8Examples.finiteCanonicalGlobalFixedPoint,BC4lean.Chapter8Examples.finiteCanonicalPointHomotopyEquiv,BC4lean.Chapter8Examples.finiteTopologicalPointSpaceEquiv,BC4lean.Chapter8Examples.reducedAssemblyHom_finitePoint_formula}
\leanok
For a finite discrete group, a lifted one-point space in the group universe
is an actual proper model. Its actual equivariant homotopy equivalence with
the canonical telescope and the compact-model colimit equivalence reduce
the canonical assembly value to the exact point cutoff/product/descent
formula in both degrees. This is a reduction of assembly, not a proof of
assembly bijectivity for every finite group.
\end{proposition}
\begin{proof}
\leanok
The canonical model's fixed-point criterion at the finite top subgroup
supplies a genuine globally fixed point. It defines the map from the lifted
point into the canonical model. The reverse map is constant; the point
composite is the identity and the other composite is equivariantly homotopic
to the identity by the actual CW mapping theorem. Apply the existing
homotopy and assembly naturality theorems and the compact-model colimit
calculation. For a nontrivial finite group, the point cutoff gives the
averaging projection; it is not replaced by the algebra unit.
\end{proof}

\begin{proposition}\label{prop:scalar-unit-projection-identity}
\uses{def:nonunital-projection-k0,prop:trivial-group-analytic-calibration}
\lean{BC4lean.KKTheory.ScalarUnitProjection.unitRelativeCycleIso,BC4lean.KKTheory.ScalarUnitProjection.relativeMap_nonunitalUnitProjection,BC4lean.KKTheory.ScalarUnitProjection.scalarKKOperatorK0AddEquiv_identity}
\leanok
The genuine scalar KK-to-operator-$K_0$ comparison sends the native scalar
identity class to the actual nonunital $K_0$ class of the unit projection.
\end{proposition}
\begin{proof}
\leanok
For the actual singleton unitization projection, the even support is the
identity and the scalar-shadow odd support is zero. The explicit coordinate
unitary identifies its whole relative-projection cycle with the standard
identity cycle, including representation, action, grading and operator.
The proved inverse comparison law then computes the actual index.
\end{proof}

\begin{proposition}\label{prop:unit-point-cutoff-hom-class}
\uses{prop:scalar-unit-projection-identity,def:cutoff-projection}
\lean{BC4lean.KKTheory.UnitPointCutoff.pointCZeroEquiv,BC4lean.KKTheory.UnitPointCutoff.reducedCutoffProjection_eq_scalarCoefficient,BC4lean.KKTheory.UnitPointCutoff.cutoffKKClass_eq_homClass}
\leanok
For the literal acting group $\mathrm{Unit}$ on an indicated compact
one-point space, the actual cutoff projection is the actual scalar
coefficient map evaluated at $1$. Its scalar KK lift is the native
homomorphism class of that coefficient map.
\end{proposition}
\begin{proof}
\leanok
Unit-group normalization gives $c(x)^2=1$. Compute the actual algebraic
projection as the identity-coordinate coefficient, integrate it into the
concrete reduced crossed product, and use genuine projection and scalar-index
naturality together with the scalar unit-projection identity. This argument
uses the literal Unit group and does not identify nontrivial finite-group
averaging with the unit.
\end{proof}

\begin{proposition}\label{prop:finite-point-cutoff-averaging}
\uses{prop:scalar-unit-projection-identity,def:cutoff-projection,prop:finite-group-point-reduction}
\lean{BC4lean.KKTheory.FinitePointCutoff.projectionScalarHom,BC4lean.KKTheory.FinitePointCutoff.cutoffKKClass_eq_homClass,BC4lean.KKTheory.FinitePointCutoff.cutoffProductHom_eq_sourcePullback,BC4lean.KKTheory.FinitePointCutoff.cutoff_apply,BC4lean.KKTheory.FinitePointCutoff.averaging,BC4lean.KKTheory.FinitePointCutoff.reducedCutoffProjection_eq_averaging,BC4lean.KKTheory.FinitePointCutoff.averaging_coe,BC4lean.KKTheory.FinitePointCutoff.cutoffKKClass_eq_averagingHomClass,BC4lean.KKTheory.FinitePointCutoff.cutoffProductHom_eq_averagingPullback}
\leanok
Let a countable discrete group $\Gamma$ act properly on a second countable
locally compact Hausdorff space $X$, and let $c$ be a normalized cutoff (one
exists exactly when the action is cocompact), with projection
$p_c$ in the reduced crossed product. Let $f\colon\mathbb C\to
C_0(X)\rtimes_r\Gamma$ be a $*$-homomorphism with $f(1)=p_c$. Then the scalar
KK lift of $[p_c]$ is the homomorphism class $[f]$, and the cutoff product
$[p_c]\otimes(-)$ is the source pullback $f^*$ in both degrees.
If $\Gamma$ is finite and $X$ is the point, every normalized cutoff is the
constant $|\Gamma|^{-1/2}$, and $p_c$ is the averaging projection
$e=\sum_{g\in\Gamma}|\Gamma|^{-1}\delta_g$. Hence the cutoff class is the
class of $z\mapsto ze$, and the cutoff product is pullback along it.
\end{proposition}
\begin{proof}
\leanok
The first part is the argument of
Proposition~\ref{prop:unit-point-cutoff-hom-class}, which does not use the
group: scalar index naturality for the native coefficient pushforward of the
scalar identity, the scalar unit-projection identity and functoriality of the
singleton projection class give $[f]\mapsto[f(1)]=[p_c]$. The cutoff product
is the left even product with $[f]$, which is the source pullback by
definition. On the point the action is trivial, so normalization reads
$|\Gamma|\,c(x)^2=1$; with $c\ge0$ this gives $c(x)=|\Gamma|^{-1/2}$, and
every coefficient $c(x)c(g^{-1}x)$ of the algebraic projection equals
$|\Gamma|^{-1}$.
\end{proof}

\begin{proposition}\label{prop:unit-scalar-descent-calibration}
\uses{def:reduced-descent,prop:descent-source-naturality}
\lean{BC4lean.KKTheory.UnitScalarDescent.evenCoefficientRestrictionIso,BC4lean.KKTheory.UnitScalarDescent.oddCoefficientRestrictionIso,BC4lean.KKTheory.UnitScalarDescent.evenClass_coefficientRestriction,BC4lean.KKTheory.UnitScalarDescent.oddClass_coefficientRestriction}
\leanok
For the literal Unit group and scalar source, pulling the actual descended
class back along the genuine scalar coefficient inclusion gives the actual
ordinary coefficient-extension class, separately in the even and odd pictures.
\end{proposition}
\begin{proof}
\leanok
The actual descended source representation restricted to the coefficient
inclusion is the actual tensor of the original representation. The identity
map on the completed coefficient-extension module intertwines every raw
cycle datum. Use the actual raw whole-cycle isomorphisms and native source
and coefficient representative theorems to pass to both full quotients.
\end{proof}

\begin{theorem}\label{thm:baum-connes-unit}
\uses{def:baum-connes-property,prop:finite-group-point-reduction,prop:unit-point-cutoff-hom-class,prop:unit-scalar-descent-calibration}
\lean{BC4lean.Chapter8Examples.unitAssemblyAddEquiv,BC4lean.Chapter8Examples.reducedAssemblyHom_unit_eq,BC4lean.Chapter8Examples.baumConnesProperty_unit}
\leanok
The actual coefficient-free reduced assembly map for the literal one-element
group $\mathrm{Unit}:\mathrm{Type}$ is an isomorphism in both degrees.
\end{theorem}
\begin{proof}
\leanok
Use the constructed canonical-to-point equivalence, the actual point source
pullback equivalence, the genuine Unit cutoff and descent identifications,
and the separately constructed even and odd scalar-to-analytic equivalences.
Prove that their composite equals Chapter~7's actual canonical assembly map.
Its explicit additive equivalence then supplies injectivity and surjectivity.
No assembly-bijectivity premise occurs. Transport to arbitrary subsingleton
groups would require an additional group-isomorphism invariance construction.
\end{proof}

\begin{proposition}\label{prop:kk-group-restriction}
\uses{def:equivariant-kk}
\lean{BC4lean.KKTheory.GroupRestriction.cycleComap,BC4lean.KKTheory.GroupRestriction.evenComap,BC4lean.KKTheory.GroupRestriction.oddComap,BC4lean.KKTheory.GroupRestriction.subsingletonParityEquiv}
\leanok
Let $f\colon \Gamma\to\Gamma'$ be a group homomorphism. Precomposing every
group action with $f$ sends $\Gamma'$-equivariant even and odd Kasparov cycles to
$\Gamma$-equivariant ones on the same module, with the same grading,
representation and operator. This respects isomorphisms and full interval
homotopies, so it defines additive restriction maps on both parity groups.
If $\Gamma$ has exactly one element and both coefficient actions are trivial,
restriction along $\mathrm{Unit}\to\Gamma$ is an additive equivalence
$KK^{\Gamma}_p(A,B)\cong KK_p(A,B)$ in both parities, with inverse the
restriction along $\Gamma\to\mathrm{Unit}$.
\end{proposition}
\begin{proof}
\leanok
Only the action field changes, and $f(g^{-1})=f(g)^{-1}$ gives the conjugation
identity; interval sections and evaluated fibres restrict definitionally. For a
one-element group both composites of the two restrictions change the action
only by $g\mapsto 1$, which is the identity, so each composite is the identity
cycle isomorphism.
\end{proof}

\section{Research and formalization status}
Keep established mathematics, completed Lean work, and open mathematical
questions separate. Completed Lean work includes the universal proper-space
construction in Chapter~4 and equivariant KK-theory, its product laws, and
the ordinary operator $K$-theory comparisons in Chapter~5, with their stated
hypotheses and documented source corrections. Chapter~6 constructs the actual
topological $K$-homology colimit and proves cofinal indexing and universal
proper-model independence under its stated hypotheses. Chapter~7 constructs the coefficient-free reduced assembly homomorphisms
in both degrees for countable discrete groups under its stated hypotheses;
a proof of the general conjecture remains a later milestone.

The verified Chapter~8 milestone consists of the exact property and conjecture
definitions, the proved model and cofinal-index comparison theorems, exact
piecewise bijectivity criteria, the geometric and analytic calibrations,
and the genuine assembly isomorphism for the literal Unit group. The general
conjecture remains unproved, so the chapter is a verified milestone rather
than a completed proof of all its targets. No numerical scalar $K$-group
calculation, nontrivial finite-group assembly isomorphism, $\mathbb Z$ case,
or $\mathrm{SL}(3,\mathbb Z)$ result is asserted by this milestone.
